\addtocontents{toc}{\protect\Needspace{3\baselineskip}}
\originalchapter{第一次测验复习：幂律宇宙中的光路}
以下题组保留 8.286 Fall 2013 Quiz 1 Review Problems 的原编号。先用统一光路积分求解，再把每道题的初值与观测量分别代入；把一条公式不加条件地套给所有题，会混淆信号相对源和相对接收者的距离。

\src{8.286 Quiz 1 Review Problem 2(a)--(b), 原题为2000年 Quiz 1 Q2；官方复习解答。}
\begin{exercise}
稳态宇宙模型虽然膨胀，却假设 $H(t)=H_0$ 和宏观物质密度 $\rho_0$ 始终不变。(a) 求最一般尺度因子 $a(t)$；(b) 求使实际密度恒定所需的单位实际体积、单位时间物质创造率。
\end{exercise}
\begin{solution}
$\dot a/a=H_0$ 给 $a(t)=A\ee^{H_0t}$，$A>0$ 是尺度规范。连续方程若加质量源 $q$，则 $\dot\rho+3H\rho=q$，故 $q=3H_0\rho_0$，单位为质量、实际体积、时间的倒数。该结果只说明\emph{稳态假说需补的源项}；普通尘埃守恒宇宙中 $q=0$，不能误认为 Friedmann 方程自发生成物质。
\end{solution}

\src{8.286 Quiz 1 Review Problem 4(a)--(d), 原题先后用于1994、2007年测验；官方复习解答。}
\begin{exercise}
平直宇宙 $a(t)=a_0\ee^{\chi t}$，$\chi>0$。在 $t=0$ 从共动原点沿正 $x$ 发光，(a) 求 $H(t)$；(b) 求光的共动轨迹 $x(t)$；(c) 若以红移 $z$ 在 $t_r$ 收到光，求 $t_r(z)$；(d) 求接收时源与接收星系的实际距离。
\end{exercise}
\begin{solution}
$H=\chi$。零线元给 $\dd x/\dd t=c/a(t)$，故 $x(t)=c[1-\ee^{-\chi t}]/(a_0\chi)$。$1+z=a(t_r)/a(0)=\ee^{\chi t_r}$，故 $t_r=\ln(1+z)/\chi$。接收时实际间距 $a(t_r)x(t_r)=(c/\chi)(\ee^{\chi t_r}-1)=cz/\chi$。此题发射刻固定为零；换其他发射时刻，应先重设光路下限。
\end{solution}

\src{8.286 Quiz 1 Review Problem 6(a)--(d), 原题为1988年 Quiz 2 Q3；官方复习解答。}
\begin{exercise}
平直模型 $a(t)=bt^{1/3}$，光于 $t_e$ 由一共动星系发出，于 $t_o$ 抵达另一共动星系。(a) 求 $t_o$ 两星系实际间距；(b) 求 $z$；(c) 将距离改写为只含 $c,t_o,z$；(d) 在 $t_e<t<t_o$ 求光到接收星系的实际距离。
\end{exercise}
\begin{solution}
共动间距 $\chi=\int_{t_e}^{t_o}(c/b)t^{-1/3}\dd t=(3c/2b)(t_o^{2/3}-t_e^{2/3})$。故
\[
D_o=\frac{3ct_o}{2}\left[1-(t_e/t_o)^{2/3}\right],\quad
1+z=(t_o/t_e)^{1/3},\quad
D_o=\frac{3ct_o}{2}[1-(1+z)^{-2}].
\]
光在 $t$ 时到接收点尚余共动路程 $(3c/2b)(t_o^{2/3}-t^{2/3})$，其实际距离为
$D_{\gamma B}(t)=(3c/2)[t^{1/3}t_o^{2/3}-t]$。
它在 $t_o$ 为零，不能把已走过的路径当作此问所求的“尚余距离”。
\end{solution}

\src{8.286 Quiz 1 Review Problem 7(a)--(e), 原题为2000年 Quiz 1 Q3；官方复习解答。}
\begin{exercise}
平直模型 $a(t)=bt^\gamma$，$\gamma>0$，今天为 $t_0$，远处共动源发光于 $t_e$、今天到达。(a) 求红移；(b) 求回看时间；(c) 以 $H_0,z,\gamma$ 求今天的实际距离；(d) 发射总功率 $P$ 各向同性时求今天通量 $J$；(e) 若源星系内部有向我们运动、相对星系速度 $v<c$ 的喷流，且喷流相对星系的短位移可略，求喷流光的观测红移 $z_J$。
\end{exercise}
\begin{solution}
$1+z=(t_0/t_e)^\gamma$、$t_0-t_e=t_0[1-(1+z)^{-1/\gamma}]$，且 $H_0=\gamma/t_0$。若 $\gamma\ne1$，光路积分给
\[
D_0=\frac{c\gamma}{H_0(1-\gamma)}
\left[1-(1+z)^{-(1-\gamma)/\gamma}\right];
\]
$\gamma=1$ 的连续极限为 $(c/H_0)\ln(1+z)$。光子能量和到达率各损失一个 $1+z$，故 $J=P/[4\pi D_0^2(1+z)^2]$。喷流局部沿光传播方向追向我们，特殊相对论源的趋近因子为 $\sqrt{(1-\beta)/(1+\beta)}$（$\beta=v/c$），因此
\[
1+z_J=(1+z)\sqrt{\frac{1-\beta}{1+\beta}}.
\]
喷流可使总红移变小甚至为蓝移；乘因子使用的是源附近局部惯性系，不把宇宙膨胀退行速度按特殊相对论直接相加。
\end{solution}

\src{8.286 Quiz 1 Review Problem 11(a)--(f), 原题为1998年 Quiz 1 Q3；官方复习解答。}
\begin{exercise}
平直模型 $a(t)=bt^{3/5}$。A 在 $t_1$ 发给 B 的光信号，当时两者实际距离 $\ell_0$。(a) 求 $H$；(b) 求收到时刻 $t_2$；(c) B 立即回信，A 于 $t_3$ 收到，求 $t_3$；(d) B 改在收到后等待任意 $\Delta t$ 回信，求 A 收到时刻 $t_4$；(e) 求迟发回信的红移；(f) 当 $\Delta t$ 很小时求 $t_4-t_3$ 的一阶项。
\end{exercise}
\begin{solution}
$H=3/(5t)$。令 $q=2/5$，则径向光路为 $\chi=(c/bq)(t_b^q-t_a^q)$，两星系固定共动距离 $\chi=\ell_0/(bt_1^{3/5})$。设 $C=(q\ell_0/c)t_1^{-3/5}$，便得到
\[
t_2^q=t_1^q+C,\quad t_3^q=t_2^q+C,\quad
t_4^q=(t_2+\Delta t)^q+C.
\]
各时刻由右式取 $1/q$ 次幂；$\Delta t$ 不要求小，故(d)不能先线性化。回信红移 $1+z=[t_4/(t_2+\Delta t)]^{3/5}$。对同一路径的相邻脉冲，接收时间微分是发射时间微分乘 $a(t_3)/a(t_2)$，故
$t_4-t_3=(t_3/t_2)^{3/5}\Delta t+O(\Delta t^2)$。
\end{solution}

\src{8.286 Quiz 1 Review Problem 9(a)--(i), 原题为1996年 Quiz 1 Q3；官方复习解答。}
\begin{exercise}
平直模型 $a(t)=bt^{3/5}$。A 在 $t_A$ 发光给 B，B 于 $t_B$ 收到。(a) 求 $H(t)$；(b) 求任意 $t$ 的实际粒子视界；(c) 求 A、B 共动间距；(d) 求两者在 $t_A$ 与 $t_B$ 的实际间距；(e) 求传播途中光同时距 A、B 相等的时刻；(f) 求 $t_A$ 时 B 相对 A 的 Hubble 退行速度；(g) 对今天的任意 $t$，求看见红移 $z$ 的源时其今天实际距离，$z\to\infty$ 极限与粒子视界比如何；(h) A 发射各向同性功率 $P$，B 收到的通量；(i) B 收光后立刻回发光，求回光抵达 A 的时刻。
\end{exercise}
\begin{solution}
置 $q=2/5$，径向光路积分是
$\chi=(c/bq)(t_B^q-t_A^q)$。
因此(a) $H=3/(5t)$；(b) $D_p=ct/q=5ct/2$；(c) 上式；
(d) $D_A=a_A\chi=(5c/2)t_A^{3/5}(t_B^{2/5}-t_A^{2/5})$，$D_B$ 同式把前因子换成 $t_B^{3/5}$；
(e) 中途到两共动端点各走半程，故 $t_{\rm mid}^{2/5}=(t_A^{2/5}+t_B^{2/5})/2$；
(f) $v_A=H_A D_A$，它是远隔共动距离变化率，可以超过 $c$。
(g) 以今天 $t$ 接收，$1+z=(t/t_e)^{3/5}$，因此
\[
D(t,z)=\frac{5ct}{2}\left[1-(1+z)^{-2/3}\right],\qquad
\lim_{z\to\infty}D(t,z)=D_p(t).
\]
(h) 今天 B 收到的通量 $J=P/[4\pi D_B^2(1+z_{AB})^2]$，其中 $1+z_{AB}=(t_B/t_A)^{3/5}$；
(i) 回程同一共动间距，若到达时为 $t_C$，则
$t_C^{2/5}-t_B^{2/5}=t_B^{2/5}-t_A^{2/5}$，即 $t_C=[2t_B^{2/5}-t_A^{2/5}]^{5/2}$。
\end{solution}

\src{8.286 Quiz 1 Review Problem 20(a)--(f), 原题为2009年 Quiz 1；官方复习解答。}
\begin{exercise}
平直尘埃宇宙 $a(t)=bt^{2/3}$。A 于 $t_1$ 向 B 发光，B 在 $t_1$ 距 A 的实际距离恰为 $c/H(t_1)$。
(a) 求信号相对 A 的实际分离速度，及发射瞬间与极晚极限；
(b) 求抵达 B 的 $t_2$；
(c) 求信号相对 B 的实际\emph{接近}速度，及两端时刻的值；
(d) 若 A 在 $t_1$ 看 B 发来的光，红移是多少？
(e) 另有 C 在 $t_1$ 也距 A 一个 Hubble 长度，且 A--C 与 A--B 垂直，C 看 B 的红移是多少？
(f) A 于 $t_1$ 各向同性发射功率 $P$，B 在 $t_2$ 收到的能流密度是多少？
\end{exercise}
\begin{solution}
径向已走实际距离
$D_{\gamma A}(t)=3c[t-t_1^{1/3}t^{2/3}]$，故分离速度
$3c-2c(t_1/t)^{1/3}$，在 $t_1$ 为 $c$，晚期趋 $3c$。
初始 B 的共动距离 $\chi_B=(3c/2b)t_1^{1/3}$。由
$3c(t_2^{1/3}-t_1^{1/3})/b=\chi_B$ 得
$t_2=(3/2)^3t_1=27t_1/8$。
信号到 B 尚余的实际距离
$D_{\gamma B}(t)=(9c/2)t_1^{1/3}t^{2/3}-3ct$，故接近速度
$-\dot D_{\gamma B}=3c[1-(t_1/t)^{1/3}]$：发射时为零，接收时为 $c$。这些是远隔\emph{实际距离}的导数；在任何局部惯性系光仍以 $c$ 传播。
若 B 的光于 $t_1$ 抵 A，发光时刻 $t_e$ 满足
$3c(t_1^{1/3}-t_e^{1/3})/b=\chi_B$，所以 $t_e=t_1/8$、$1+z_{BA}=(t_1/t_e)^{2/3}=4$，即 $z_{BA}=3$。B--C 在平直空间的共动间距是 $\sqrt2\chi_B$，故
$t_{e,BC}^{1/3}=t_1^{1/3}(1-1/\sqrt2)$，
$z_{BC}=(1-1/\sqrt2)^{-2}-1$。
最后 A--B 在 $t_2$ 的实际间距为 $D_B(t_2)=ct_2$，$1+z_{AB}=(t_2/t_1)^{2/3}=9/4$，故
\[
J=\frac{P}{4\pi[(9/4)ct_2]^2}
=\frac{P}{4\pi[(243/32)ct_1]^2}.
\]
衰减因子含波面面积及能量、到达率各一重红移。
\end{solution}

\src{8.286 Quiz 1 Review Problem 12 与 Problem 13；官方复习解答。}
\begin{exercise}
(12) 在仅有无压尘埃、可有空间曲率的宇宙，减速参数 $q=-a\ddot a/\dot a^2$ 与 $\Omega=\rho/\rho_c$ 有何关系？
(13) 平直、仅光子辐射主导且 $\rho\propto a^{-4}$ 的宇宙，求大爆炸后 $a(t)$ 的幂律。
\end{exercise}
\begin{solution}
尘埃加速度式 $\ddot a/a=-4\pi G\rho/3$，除以 $H^2$ 得 $q=4\pi G\rho/(3H^2)=\Omega/2$。真空能或其他压强若存在则关系需增项，不能只看总 $\Omega$。辐射题用 $H^2=(8\pi G/3)\rho\propto a^{-4}$，所以 $\dot a\propto a^{-1}$、$a\propto t^{1/2}$；积分常数以大爆炸点 $a(0)=0$ 固定。
\end{solution}

\src{8.286 Quiz 1 Review Problem 17(a)--(e), 原题为2005年 Quiz 1 Q3；官方复习解答。}
\begin{exercise}
在平直辐射宇宙 $a(t)=bt^{1/2}$，一光子今天 $t_f$ 到原点，其发射源今天恰在粒子视界上。(a) 求今天实际粒子视界；(b) 求发射时刻；(c) 求源的共动距离；(d) 光子在 $0\le t\le t_f$ 到原点的实际距离；(e) 其最大实际距离与时刻。
\end{exercise}
\begin{solution}
粒子视界 $D_p(t_f)=a_f\int_0^{t_f}c\,\dd t/a=2ct_f$。恰在视界的极限光线从 $t_e\to0^+$ 发出，共动源距离 $\chi=2c\sqrt{t_f}/b$。光子的剩余共动距离为 $(2c/b)(\sqrt{t_f}-\sqrt t)$，故
\[
D_{\gamma O}(t)=2c(\sqrt{tt_f}-t).
\]
对 $t$ 求导，$c\sqrt{t_f/t}-2c=0$ 给 $t_{\max}=t_f/4$、$D_{\max}=ct_f/2$。发射时实际距离趋零，但共动距离非零；“光从视界发出”是 $t_e\to0$ 极限，不要求物理源存在于奇点。
\end{solution}

\src{8.286 Quiz 1 Review Problem 22(a)--(e), 原题为2011年 Quiz 1；官方复习解答。}
\begin{exercise}
平直尘埃宇宙 $a(t)=bt^{2/3}$，今天 $t_0$，来自今日视界上的共动源的光正到原点。(a) 求今天实际视界；(b) 求 $t=0$ 光的共动坐标 $x_0$；(c) 求 $x(t)$；(d) 求光到原点的实际距离；(e) 求途中该距离的最大值和时刻。
\end{exercise}
\begin{solution}
$D_p(t_0)=3ct_0$，所以 $x_0=D_p/a_0=3ct_0^{1/3}/b$。向原点的径向光路满足 $\dot x=-c/a$，积分给
\[
x(t)=\frac{3c}{b}(t_0^{1/3}-t^{1/3}),\quad
D_{\gamma O}(t)=3c(t_0^{1/3}t^{2/3}-t).
\]
对后式求导为零得 $t_{\max}=8t_0/27$，再代回得 $D_{\max}=4ct_0/9$。$D_{\gamma O}$ 最初可增大，因为空间膨胀率大于光靠近的局部速度；最终光仍抵达原点。这是实际距离曲线，不与始终单调减小的共动剩余距离混同。
\end{solution}

\src{8.286 Quiz 1 Review Problem 21(a)--(d), 原题为2011年 Quiz 1；官方复习解答署名 Daniele Bertolini。}
\begin{exercise}
(a) Hubble 当年用何种办法估计星系距离并建立距离与红移关系？(b) 何为宇宙学原理？(c) 区分均匀与各向同性，并在欧氏空间判断两句：仅关于一个点各向同性必均匀；关于两个不同点各向同性必均匀。(d) 若允许曲面，关于两个点各向同性还必均匀吗？给例子。
\end{exercise}
\begin{solution}
(a) 标准烛光方法：Cepheid 脉动周期标定本征光度，测其视亮度求距离；1929年更远的约18个星系还借最亮恒星的视亮度作为距离阶梯，不能说全部距离都由其星系内 Cepheid 直接测得。(b) 在足够大的尺度上，典型共动观测者所见的统计宇宙没有特殊地点与方向。(c) 均匀指无特殊位置，各向同性指无特殊方向。只有一中心各向同性不够，例如密度 $\rho=\rho(|\boldsymbol x|)$ 非常数。欧氏空间两个不同中心 A、B 都各向同性时，对于确定性的标量密度，$\rho=F(|x-A|)=G(|x-B|)$；固定 $|x-A|=r$ 的圆或球上 $|x-B|$ 扫过一个区间，$G$ 在该区间相同，相邻半径的区间互相重叠，推出全空间常数。一般“典型观察”也需将各向同性定义为\emph{全部统计性质}的旋转不变，而非仅一项观测。(d) 曲面反例是球面上依纬度而变的密度：北极、南极各自看各方向相同，却并非处处均匀。官方原解指出书中“两圆相交”的简证对所有点不充分；这里的重叠区间论证给出欧氏标量情形所需的连通步骤。
\end{solution}
